\documentclass{article}
\usepackage{amsmath, amssymb, amsthm}
\usepackage{hyperref}
\usepackage{geometry}
\geometry{margin=1in}

\title{Erd\H{o}s Problem \#893}
\date{Last updated: 2025-08-31}
\author{Source: erdosproblems.com}

\begin{document}
\maketitle

\noindent\textbf{Status:} OPEN \quad \textbf{No prize}

\noindent\textbf{Tags:} number theory, divisors

\medskip
\noindent\textbf{Problem Statement:}

If $\tau(n)$ counts the divisors of $n$ then let\[f(n)=\sum_{1\leq k\leq n}\tau(2^k-1).\]Does $f(2n)/f(n)$ tend to a limit?

\bigskip
\noindent\textbf{Remarks:}

Erd\H{o}s \cite{Er98} says that 'probably there is no simple asymptotic formula for $f(n)$ since $f(n)$ increases too fast'. 

Kova\v{c} and Luca \cite{KoLu25} (building on a heuristic independently found by Cambie (personal communication)) have shown that there is no finite limit, in that\[\limsup_{n\to \infty}\frac{f(2n)}{f(n)}=\infty,\]and provide both theoretical and numerical evidence that suggests $\lim \frac{f(2n)}{f(n)}=\infty$.

\bigskip
\noindent\textbf{References:}

[Er98] Erd\H{o}s, Paul, Some of my new and almost new problems and results in combinatorial number theory. Number theory (Eger, 1996) (1998), 169-180.
 
 [KoLu25] V. Kova\v{c} and F. Luca, On the number of divisors of Mersenne numbers. arXiv:2506.04883 (2025).

\bigskip
\noindent\small{Source: \url{https://www.erdosproblems.com/893}}
\end{document}
